\let\R\relax
\newcommand{\R}{\mathbb{R}}
\let\C\relax
\newcommand{\C}{\mathbb{C}}
\let\N\relax
\newcommand{\N}{\mathbb{N}}
\let\Z\relax
\newcommand{\Z}{\mathbb{Z}}
\let\Q\relax
\newcommand{\Q}{\mathbb{Q}}
\let\ch\relax
\newcommand{\ch}{\operatorname{ch}}
\let\Td\relax
\newcommand{\Td}{\operatorname{Td}}
\let\sign\relax
\newcommand{\sign}{\operatorname{sign}}
\let\cL\relax
\newcommand{\cL}{\mathcal{L}}

\chapter{Friedrich Hirzebruch (1988)}
\index{Hirzebruch, Friedrich}

\begin{quote}
\it ``Hirzebruch's work is a model of elegance and depth. His
Hirzebruch--Riemann--Roch theorem and the Hirzebruch signature theorem
are cornerstones of modern mathematics, linking analysis, topology, and
algebraic geometry. His creation of the Max Planck Institute for
Mathematics in Bonn established a world center for mathematical
research.'' --- Wolf Prize Citation
\end{quote}

Friedrich Ernst Peter Hirzebruch (1927--2012) was one of the most
influential mathematicians of the postwar era. His work bridges
topology, algebraic geometry, and differential geometry with
extraordinary clarity and power. The 1988 Wolf Prize in Mathematics
(shared with Lars H\"ormander) recognized his fundamental contributions:
the Hirzebruch--Riemann--Roch theorem, the Hirzebruch signature theorem,
and his systematic development of multiplicative sequences and
characteristic classes in algebraic geometry.

Hirzebruch's mathematics is characterized by a rare synthesis of
topological and algebraic methods. He transformed the classical
Riemann--Roch problem\index{Riemann-Roch theorem} from a theorem about curves into a general
principle for complex manifolds, revealing the deep connection between
the analysis of sheaf cohomology, the topology of characteristic
classes, and the geometry of algebraic varieties. His 1956 book
\textit{Neue topologische Methoden in der algebraischen Geometrie}
became an instant classic and shaped the development of modern algebraic
geometry.

\section{Main Contributions}

\begin{itemize}
    \item \textbf{Hirzebruch--Riemann--Roch Theorem (1954):}
    The generalization of the classical Riemann--Roch theorem from
    algebraic curves to arbitrary compact complex manifolds. For a
    holomorphic vector bundle \(E\) over a compact complex manifold
    \(X\), the Euler characteristic \(\chi(X, E)\) is expressed as
    \[
    \chi(X, E) = \int_X \ch(E) \Td(TX),
    \]
    where \(\ch(E)\) is the Chern character of \(E\) and \(\Td(TX)\) is
    the Todd class of the tangent bundle. This theorem unified sheaf
    cohomology with characteristic class theory and paved the way for
    the Atiyah--Singer index theorem.

    \item \textbf{Hirzebruch Signature Theorem (1956):}\index{Hirzebruch signature theorem}
    For a compact oriented smooth manifold \(M\) of dimension \(4k\),
    the signature \(\tau(M)\) equals a certain polynomial in the
    Pontryagin classes\index{Pontryagin class}, expressed through the \(L\)-genus:
    \[
    \tau(M) = \int_M L(p_1, \dots, p_k).
    \]
    For a \(4\)-manifold this simplifies to
    \(\tau(M) = \frac13 \int_M p_1(M)\). The signature theorem is one
    of the most striking results connecting topology and analysis.

    \item \textbf{Multiplicative Sequences:}
    Hirzebruch developed the theory of multiplicative sequences of
    polynomials, which associates to every power series
    \(Q(x) = 1 + q_1 x + q_2 x^2 + \cdots\) a sequence of polynomials
    \(K_n(p_1, \dots, p_n)\) such that the genus
    \(\int_M K(p_1, \dots, p_k)\) is a multiplicative invariant. The
    Todd genus arises from \(Q(x) = x/(1 - e^{-x})\), the \(L\)-genus
    from \(Q(x) = \sqrt{x}/\tanh\sqrt{x}\), and the \(\hat{A}\)-genus
    from \(Q(x) = \sqrt{x}/ \sinh\sqrt{x}\).

    \item \textbf{Hirzebruch Surfaces:}
    The ruled surfaces \(\Sigma_n = \mathbb{P}(\mathcal{O} \oplus
    \mathcal{O}(-n))\) over \(\mathbb{P}^1\), now called Hirzebruch
    surfaces. They provide fundamental examples in the classification
    of algebraic surfaces, with \(\Sigma_0 = \mathbb{P}^1 \times
    \mathbb{P}^1\) and \(\Sigma_1\) the blowup of \(\mathbb{P}^2\) at a
    point.

    \item \textbf{Topology of Algebraic Varieties:}
    Systematic application of characteristic classes to algebraic
    varieties, including the computation of Chern numbers, the study
    of cobordism classes, and the development of Chern numbers for
    algebraic surfaces.

    \item \textbf{Max Planck Institute for Mathematics:}
    In 1981, Hirzebruch founded the Max Planck Institute for
    Mathematics in Bonn, which became a world center for research. The
    institute's Hirzebruch Prize (established 2018) honors his legacy.

    \item \textbf{Proportionality Principle and Index Theory:}
    His work on the proportionality principle in complex geometry, the
    relation between Chern and Pontryagin numbers, and the topology of
    complex flag manifolds. His collaboration with Atiyah and others
    on the index of elliptic operators and \(K\)-theory.
\end{itemize}

\section{Origin of the Problem}

\subsection{The Classical Riemann--Roch Theorem}

The Riemann--Roch theorem for algebraic curves is one of the most
beautiful results of nineteenth-century mathematics. For a compact
Riemann surface \(C\) of genus \(g\) and a divisor \(D\) on \(C\), the
theorem states
\[
\ell(D) - \ell(K - D) = \deg D + 1 - g,
\]
where \(\ell(D)\) is the dimension of the space of meromorphic functions
with poles bounded by \(D\) and \(K\) is the canonical divisor. In
sheaf-theoretic language, this becomes
\[
\chi(C, \mathcal{O}(D)) = \deg D + 1 - g,
\]
where
\(\chi(C, \mathcal{O}(D)) = \dim H^0(C, \mathcal{O}(D)) -
\dim H^1(C, \mathcal{O}(D))\).

The classical theorem was a powerful tool, but it was limited to curves.
Mathematicians in the first half of the twentieth century sought
generalizations to higher-dimensional algebraic varieties. The problem
was: given a line bundle (or vector bundle) on a smooth projective
variety, can one compute the Euler characteristic of its sheaf
cohomology in terms of topological invariants?

The difficulty lay in the fact that higher-dimensional varieties have
many intermediate cohomology groups, and the alternating sum
\[
\chi(X, E) = \sum_{i=0}^{\dim X} (-1)^i \dim H^i(X, \mathcal{O}(E))
\]
must be expressed in terms of characteristic classes that capture the
topology of both \(E\) and the tangent bundle \(TX\).

\subsection{The Signature Problem}

Parallel to the Riemann--Roch problem was the question of the signature
of a manifold. For a compact oriented \(4k\)-dimensional manifold
\(M\), the intersection form on the middle-dimensional cohomology
\[
Q_M: H^{2k}(M, \R) \times H^{2k}(M, \R) \to \R,
\quad
Q_M(\alpha, \beta) = \int_M \alpha \cup \beta
\]
is a symmetric bilinear form. Its signature \(\tau(M)\) is the number of
positive eigenvalues minus the number of negative eigenvalues.

It was observed that the signature behaves multiplicatively under
coverings and satisfies a certain additivity property under cutting and
pasting. These properties suggested that the signature should be
expressible as a linear combination of Pontryagin numbers --- integrals
of products of Pontryagin classes evaluated on the fundamental cycle.
But what combination? Euler's theorem on the Euler characteristic as a
sum of Betti numbers had a topological formula; the signature cried out
for a similar expression.

\subsection{Multiplicative Sequences and Genera}

Hirzebruch's key insight came from the theory of multiplicative
sequences. Suppose we want to construct an invariant \(I(M)\) of
oriented closed manifolds that is:
\begin{enumerate}
    \item \textbf{Multiplicative:}
    \(I(M_1 \times M_2) = I(M_1) \cdot I(M_2)\),
    \item \textbf{Additive:}
    \(I(M_1 \sqcup M_2) = I(M_1) + I(M_2)\),
    \item \textbf{Cobordism invariant:}
    if \(M = \partial W\), then \(I(M) = 0\).
\end{enumerate}

Such an invariant is called a \textit{genus}. Hirzebruch recognized that
genera correspond precisely to multiplicative sequences of polynomials,
and that each genus is determined by a power series. The key examples
are:
\begin{itemize}
    \item The \textbf{Todd genus} \(\Td(X)\) of a complex manifold
    \(X\), corresponding to \(Q(x) = x/(1 - e^{-x})\).
    \item The \(L\)-{\bf genus} \(L(M)\) of an oriented manifold,
    corresponding to \(Q(x) = \sqrt{x}/\tanh\sqrt{x}\).
    \item The \(\hat{A}\)-{\bf genus} of a spin manifold, corresponding
    to \(Q(x) = \sqrt{x}/\sinh\sqrt{x}\).
\end{itemize}

The signature theorem would be \(\tau(M) = L(M)[M]\), and the HRR
theorem would be \(\chi(X, \mathcal{O}) = \Td(TX)[X]\) (the Todd genus).

\section{How It Was Solved}

\subsection{The Hirzebruch--Riemann--Roch Theorem (1954)}

Hirzebruch's breakthrough in 1954 was to realize that the Chern
character \(\ch(E)\) and the Todd class \(\Td(TX)\) are the correct
cohomological ingredients. The Chern character of a vector bundle
\(E\) is defined from the Chern classes via the formal factorization:
if \(E\) splits formally into line bundles with Chern roots \(x_i\),
then
\[
\ch(E) = \sum_{i=1}^{\rank E} e^{x_i}
= \rank E + c_1(E) + \frac12(c_1^2 - 2c_2)(E)
+ \frac16(c_1^3 - 3c_1c_2 + 3c_3)(E) + \cdots.
\]

The Todd class of the tangent bundle \(TX\) is defined via the
multiplicative sequence associated to \(Q(x) = x/(1 - e^{-x})\): if the
Chern roots of \(TX\) are \(y_i\), then
\[
\Td(TX) = \prod_{i=1}^n \frac{y_i}{1 - e^{-y_i}}
= 1 + \frac12 c_1 + \frac1{12}(c_1^2 + c_2) + \cdots.
\]

\begin{theorem}[Hirzebruch--Riemann--Roch, 1954]
Let \(X\) be a compact complex manifold of dimension \(n\) and let
\(E\) be a holomorphic vector bundle on \(X\). Then
\[
\chi(X, E) = \int_X \ch(E) \Td(TX),
\]
where \(\chi(X, E) = \sum_{i=0}^n (-1)^i \dim_{\C} H^i(X,
\mathcal{O}(E))\) is the Euler characteristic, \(\ch(E)\) is the Chern
character of \(E\), and \(\Td(TX)\) is the Todd class of the tangent
bundle \(TX\). The integral denotes evaluation of the degree-\(2n\)
component of the cohomology class on the fundamental cycle of \(X\).
\end{theorem}

\begin{example}[Riemann--Roch for Curves]
For a compact Riemann surface \(C\) of genus \(g\) and a line bundle
\(L = \mathcal{O}(D)\) of degree \(d\), we have
\[
\ch(L) = 1 + c_1(L) = 1 + d\,\omega,
\]
where \(\omega\) is the generator of \(H^2(C, \Z)\). The Todd class of
\(TC\) is
\[
\Td(TC) = 1 + \frac12 c_1(TC) = 1 + (1 - g)\omega.
\]
Then
\[
\chi(C, L) = \int_C (1 + d\omega)(1 + (1-g)\omega)
= \int_C (d + 1 - g)\omega = d + 1 - g,
\]
which is exactly the classical Riemann--Roch theorem.
\end{example}

\begin{example}[Projective Space]
For \(X = \mathbb{P}^n\) with the hyperplane class
\(h \in H^2(\mathbb{P}^n, \Z)\) and line bundle \(\mathcal{O}(d)\),
we have
\[
\ch(\mathcal{O}(d)) = e^{dh},
\qquad
\Td(T\mathbb{P}^n) = \left(\frac{h}{1 - e^{-h}}\right)^{\!n+1}.
\]
The Euler characteristic is the coefficient of \(h^n\) in
\(e^{dh} (h/(1 - e^{-h}))^{n+1}\), giving
\[
\chi(\mathbb{P}^n, \mathcal{O}(d)) = \binom{d + n}{n}
\]
for \(d \ge 0\).
\end{example}

Hirzebruch proved this theorem by using the splitting principle and the
theory of multiplicative sequences. He reduced the general case to the
case of line bundles over projective space by constructing a suitable
embedding and using the fact that the Chern character and Todd class
behave naturally under pushforward.

The HRR theorem was a spectacular advance. It generalized a 1953 result
of Kodaira and Spencer and immediately yielded many new computations.
It also laid the groundwork for the Grothendieck--Riemann--Roch theorem
(1958), which generalized HRR to proper morphisms of schemes, and for
the Atiyah--Singer index theorem (1963), which unified HRR and the
signature theorem under the umbrella of elliptic operators.

\subsection{The Signature Theorem (1956)}

The Hirzebruch signature theorem expresses the signature of a manifold
in terms of its Pontryagin classes. For an oriented manifold \(M\), the
Pontryagin classes \(p_i(M) \in H^{4i}(M, \Z)\) are characteristic
classes that capture the topology of the tangent bundle.

\begin{theorem}[Hirzebruch Signature Theorem, 1956]
Let \(M\) be a compact oriented smooth manifold of dimension \(4k\).
The signature \(\tau(M)\) is given by
\[
\tau(M) = \int_M L_k(p_1, \dots, p_k),
\]
where \(L_k\) is the \(k\)-th polynomial in the multiplicative sequence
associated to the power series
\[
Q(x) = \frac{\sqrt{x}}{\tanh\sqrt{x}}
= 1 + \frac13 x - \frac1{45} x^2 + \cdots.
\]
The first few \(L\)-polynomials are:
\[
\begin{aligned}
L_1 &= \frac13 p_1,\\[2pt]
L_2 &= \frac1{45}(7p_2 - p_1^2),\\[2pt]
L_3 &= \frac1{945}(62p_3 - 13p_1 p_2 + 2p_1^3),\\[2pt]
L_4 &= \frac1{14175}(381p_4 - 71p_1 p_3 - 19p_2^2 + 22p_1^2 p_2
- 3p_1^4).
\end{aligned}
\]
For a \(4\)-manifold, this reduces to
\[
\tau(M) = \frac13 \int_M p_1(M).
\]
\end{theorem}

\begin{example}[Signature of Complex Projective Space]
The complex projective plane \(\mathbb{CP}^2\) has
\(p_1(\mathbb{CP}^2) = 3h^2\) where \(h\) generates
\(H^2(\mathbb{CP}^2, \Z)\). Then
\[
\tau(\mathbb{CP}^2) = \frac13 \int_{\mathbb{CP}^2} 3h^2 = 1.
\]
The intersection form on \(H^2(\mathbb{CP}^2, \R)\) is positive definite
of rank 1, so \(\tau = 1\) as expected.
\end{example}

\begin{example}[Signature of a K3 Surface]
A K3 surface is a compact complex surface with \(c_1 = 0\). Its
signature is \(\tau = -16\), and the signature theorem gives
\[
-16 = \frac13 \int p_1,
\]
relating the signature, Euler characteristic \(e = 24\), and the first
Pontryagin class.
\end{example}

The signature theorem was proven by applying the theory of
multiplicative sequences to the Pontryagin classes. Hirzebruch used the
fact that the signature is a genus --- a multiplicative, additive,
cobordism-invariant function on oriented manifolds --- and showed that
any genus must be a linear combination of Pontryagin numbers, with
coefficients determined by the corresponding multiplicative sequence.

\subsection{Multiplicative Sequences}

The theory of multiplicative sequences is one of Hirzebruch's most
enduring contributions. It provides a systematic method for
constructing characteristic numbers.

\begin{definition}[Multiplicative Sequence]
Let \(Q(x) = 1 + q_1 x + q_2 x^2 + \cdots\) be a formal power series
with \(q_k \in \Q\). The associated \textbf{multiplicative sequence}
\(\{K_n\}\) is defined by the formal factorization
\[
\sum_{n=0}^\infty K_n(p_1, \dots, p_n) t^n
= \prod_{i=1}^\infty Q(x_i t),
\]
where \(p_k\) is the \(k\)-th elementary symmetric function in the
variables \(x_i\). The polynomials \(K_n\) are homogeneous of weight
\(n\) (where \(p_k\) has weight \(k\)).
\end{definition}

\begin{theorem}[Hirzebruch, 1956]
Every genus (multiplicative, additive, cobordism-invariant function) on
oriented manifolds is given by evaluating an \(L\)-polynomial on the
fundamental class. On complex manifolds, every multiplicative invariant
is given by evaluating a Todd-type polynomial on the Chern classes. The
correspondence between power series and genera is bijective:
\[
Q(x) \longleftrightarrow \text{genus } \text{gen}_Q(M)
= \int_M K_Q(p_1, \dots, p_k).
\]
\end{theorem}

The key examples are:

\[
\begin{array}{l|l|l}
\textbf{Genus} & \textbf{Power series } Q(x) & \textbf{Application}
\\ \hline
\text{Todd genus} & x/(1 - e^{-x}) & \text{Complex manifolds, HRR}
\\
L\text{-genus} & \sqrt{x}/\tanh\sqrt{x} & \text{Signature theorem}
\\
\hat{A}\text{-genus} & \sqrt{x}/\sinh\sqrt{x} & \text{Spin manifolds}
\\
\text{Euler characteristic} & x/\tanh x & \text{Euler class}
\end{array}
\]

\subsection{Hirzebruch Surfaces}

The Hirzebruch surfaces, introduced by Hirzebruch in the 1950s, are a
family of complex surfaces that serve as fundamental examples in
algebraic geometry.

\begin{definition}[Hirzebruch Surface]
For \(n \ge 0\), the Hirzebruch surface \(\Sigma_n\) is the projective
bundle
\[
\Sigma_n = \mathbb{P}\bigl(\mathcal{O}_{\mathbb{P}^1}
\oplus \mathcal{O}_{\mathbb{P}^1}(-n)\bigr)
\]
over \(\mathbb{P}^1\). Equivalently, it is the total space of a
\(\mathbb{P}^1\)-bundle over \(\mathbb{P}^1\).
\end{definition}

These surfaces have the following properties:
\begin{itemize}
    \item \(\Sigma_0 = \mathbb{P}^1 \times \mathbb{P}^1\), the quadric
    surface.
    \item \(\Sigma_1\) is isomorphic to the blowup of \(\mathbb{P}^2\)
    at a point.
    \item For \(n \ge 2\), \(\Sigma_n\) is not isomorphic to a blowup
    of \(\mathbb{P}^2\) but is a ruled surface with a unique section
    of self-intersection \(-n\).
    \item All \(\Sigma_n\) are rational surfaces, diffeomorphic to each
    other for fixed parity of \(n\).
    \item The Chern numbers are \(c_1^2(\Sigma_n) = 8\) and
    \(c_2(\Sigma_n) = 4\) for all \(n\).
\end{itemize}

Hirzebruch surfaces play an important role in:
\begin{enumerate}
    \item \textbf{Classification of algebraic surfaces:} They are the
    minimal rational surfaces (together with \(\mathbb{P}^2\)).
    \item \textbf{The geography of surfaces:} They lie on the Noether
    line \(c_1^2 = 2c_2 - 4\).
    \item \textbf{Toric geometry:} Each \(\Sigma_n\) is a toric variety.
    \item \textbf{Symplectic geometry:} They provide fundamental
    examples in the study of symplectic 4-manifolds.
\end{enumerate}

\section{Mathematical Theory}

\subsection{The Chern Character}

Let \(E\) be a complex vector bundle of rank \(r\) over a smooth
manifold \(X\). The total Chern class is
\[
c(E) = 1 + c_1(E) + c_2(E) + \cdots + c_r(E)
\in H^{\bullet}(X, \Z).
\]
The splitting principle allows us to write
\(c(E) = \prod_{i=1}^r (1 + x_i)\) where the formal Chern roots \(x_i\)
lie in \(H^2(X, \Z)\) after passing to a suitable flag manifold.

\begin{definition}[Chern Character]
The \textbf{Chern character} of \(E\) is
\[
\ch(E) = \sum_{i=1}^r e^{x_i}
= r + \sum_{i=1}^r x_i + \frac12 \sum_{i=1}^r x_i^2 + \cdots
\in H^{\bullet}(X, \Q).
\]
In terms of Chern classes, the first few terms are:
\[
\begin{aligned}
\ch_0(E) &= r,\\
\ch_1(E) &= c_1(E),\\
\ch_2(E) &= \tfrac12\bigl(c_1(E)^2 - 2c_2(E)\bigr),\\
\ch_3(E) &= \tfrac16\bigl(c_1(E)^3 - 3c_1(E)c_2(E)
+ 3c_3(E)\bigr).
\end{aligned}
\]
\end{definition}

The Chern character is a ring homomorphism from \(K\)-theory to
cohomology:
\[
\ch(E \oplus F) = \ch(E) + \ch(F),\qquad
\ch(E \otimes F) = \ch(E)\,\ch(F).
\]

\subsection{The Todd Class}

\begin{definition}[Todd Class]
The \textbf{Todd class} of a complex vector bundle \(E\) with Chern
roots \(x_i\) is
\[
\Td(E) = \prod_{i=1}^r \frac{x_i}{1 - e^{-x_i}}
\in H^{\bullet}(X, \Q).
\]
The first few terms in terms of Chern classes are:
\[
\begin{aligned}
\Td_0 &= 1,\\
\Td_1 &= \tfrac12 c_1,\\
\Td_2 &= \tfrac1{12}(c_1^2 + c_2),\\
\Td_3 &= \tfrac1{24} c_1 c_2,\\
\Td_4 &= \tfrac1{720}(-c_1^4 + 4c_1^2 c_2 + 3c_2^2
+ c_1 c_3 - c_4).
\end{aligned}
\]
\end{definition}

The Todd class is multiplicative:
\(\Td(E \oplus F) = \Td(E) \Td(F)\). For a complex manifold \(X\), the
Todd genus \(\Td(X) = \int_X \Td(TX)\) equals the arithmetic genus
\(\chi(X, \mathcal{O}_X)\) by the HRR theorem.

\subsection{The Hirzebruch--Riemann--Roch Theorem in Detail}

For a compact complex manifold \(X\) of dimension \(n\) and a
holomorphic vector bundle \(E\), the Euler characteristic can be
expanded in terms of the degree-\(2n\) component:
\[
\chi(X, E) = \int_X \bigl[\ch(E) \Td(TX)\bigr]_{2n}.
\]

\begin{example}[Chern Numbers of a Surface]
For a complex surface \(S\) with Chern classes \(c_1, c_2\) and a line
bundle \(L\) with \(c_1(L) = D\), the HRR theorem gives
\[
\chi(S, L) = \frac12 D^2 + \frac12 D \cdot c_1
+ \frac1{12}(c_1^2 + c_2).
\]
When \(L = \mathcal{O}_S\), we recover the Noether formula
\[
\chi(S, \mathcal{O}_S) = \frac1{12}(c_1^2 + c_2).
\]
\end{example}

The HRR theorem was generalized by Grothendieck in 1958 to the
Grothendieck--Riemann--Roch (GRR) theorem for proper morphisms
\(f: X \to Y\):
\[
\ch(f_! E) \Td(TY) = f_*(\ch(E) \Td(TX)),
\]
where \(f_!\) denotes the direct image in \(K\)-theory and \(f_*\) the
pushforward in cohomology. The GRR theorem is a central tool in modern
algebraic geometry, used to study moduli spaces, intersection theory,
and the geometry of families.

\subsection{The Signature Theorem and the \texorpdfstring{\(L\)}{L}-Genus}

The \(L\)-genus is defined via the multiplicative sequence associated
to \(Q(x) = \sqrt{x}/\tanh\sqrt{x}\). Its expansion is
\[
\frac{\sqrt{x}}{\tanh\sqrt{x}}
= 1 + \sum_{k=1}^\infty \frac{2^{2k} B_{2k}}{(2k)!} x^k,
\]
where \(B_{2k}\) are Bernoulli numbers. The first Bernoulli numbers are
\(B_2 = 1/6\), \(B_4 = -1/30\), \(B_6 = 1/42\), giving the coefficients
\(1/3, -1/45, 2/945, \dots\) in the \(L\)-polynomials.

\begin{theorem}[Signature Theorem, Equivalent Form]
For a compact oriented \(4k\)-manifold \(M\),
\[
\tau(M) = \int_M L_k(p_1, \dots, p_k)
= L_k(p_1, \dots, p_k)[M].
\]
In particular, the signature is a rational linear combination of
Pontryagin numbers.
\end{theorem}

The signature theorem has important consequences:
\begin{itemize}
    \item The signature of a \(4k\)-dimensional manifold is an integer,
    so the Pontryagin numbers must satisfy certain integrality
    conditions --- the famous Hirzebruch integrality theorems.
    \item For a \(4\)-manifold, \(\int_M p_1(M)\) must be divisible by
    \(3\).
    \item The theorem implies that the signature is a cobordism
    invariant, and the \(L\)-genus is a ring homomorphism from the
    oriented cobordism ring to \(\Q\).
\end{itemize}

\subsection{The Proportionality Principle}

One of Hirzebruch's deep insights was the proportionality principle:
for compact quotients of symmetric spaces, the characteristic numbers
are proportional to those of the corresponding compact dual symmetric
space.

Let \(G/K\) be a Hermitian symmetric space of compact type and
\(\Gamma \backslash G/K\) a compact quotient of the non-compact dual.
Then Hirzebruch showed
\[
\frac{c_n(\Gamma \backslash G/K)}{\operatorname{Vol}(\Gamma \backslash G/K)}
= \frac{c_n(G/K)}{\operatorname{Vol}(G/K)},
\]
where \(c_n\) denotes any Chern number. This principle connects the
topology of locally symmetric spaces to their volumes and has deep
applications in the theory of automorphic forms.

\subsection{Connections to the Atiyah--Singer Index Theorem}

The Atiyah--Singer index theorem (1963) unified the HRR theorem and the
signature theorem as special cases of a single principle: the index of
an elliptic operator on a compact manifold equals a topological
invariant expressed in terms of characteristic classes.

Specifically:
\begin{itemize}
    \item The HRR theorem is the index theorem for the
    \(\bar\partial\)-operator on a complex manifold:
    \(\operatorname{Ind}(\bar\partial_E) = \chi(X, E)
    = \int_X \ch(E) \Td(TX)\).
    \item The signature theorem is the index theorem for the signature
    operator \(d + d^*\) on a Riemannian manifold:
    \(\operatorname{Ind}(\text{signature operator}) = \tau(M)
    = \int_M L(p_1, \dots, p_k)\).
    \item The \(\hat{A}\)-genus appears as the index of the Dirac
    operator on a spin manifold.
\end{itemize}

Hirzebruch's work provided the essential blueprint for the index
theorem, and Atiyah and Singer explicitly credit Hirzebruch's
formulations as the inspiration for their general theory.

\subsection{Hirzebruch Surfaces: Further Details}

The Hirzebruch surface \(\Sigma_n\) can be described explicitly. Let
\(H\) be the divisor class of a fiber of the projection
\(\pi: \Sigma_n \to \mathbb{P}^1\) and let \(E\) be the class of the
section with self-intersection \(-n\) (the exceptional section). Then
the intersection pairing is
\[
E^2 = -n,\quad E \cdot H = 1,\quad H^2 = 0.
\]
The canonical divisor is \(K_{\Sigma_n} = -2E - (n+2)H\), and the Chern
numbers are
\[
c_1(\Sigma_n)^2 = 8,\quad c_2(\Sigma_n) = 4.
\]

The Hirzebruch surfaces have been generalized to higher dimensions: the
so-called Hirzebruch--Grothendieck manifolds
\(\mathbb{P}(\mathcal{O} \oplus \mathcal{O}(-n_1)
\oplus \cdots \oplus \mathcal{O}(-n_k))\) over \(\mathbb{P}^1\), and
more generally projective bundles over arbitrary bases.

\section{Impact on Science and Technology}

\subsection{Algebraic Geometry}

The Hirzebruch--Riemann--Roch theorem transformed algebraic geometry.
Before 1954, computations of sheaf cohomology on higher-dimensional
varieties were extremely difficult. After HRR, the Euler characteristic
of any vector bundle could be computed mechanically from its Chern
classes.

The theorem was essential for:
\begin{itemize}
    \item \textbf{Classification of algebraic surfaces:} Kodaira's
    classification used HRR to compute invariants.
    \item \textbf{Moduli spaces:} The dimension and properties of
    moduli spaces of curves and vector bundles depend on HRR-type
    computations. For example, the dimension of the moduli space
    \(\mathcal{M}_g\) of curves of genus \(g\) is \(3g-3\), computed by
    Riemann--Roch.
    \item \textbf{Intersection theory:} HRR is the foundation for the
    Riemann--Roch theorem on algebraic varieties, generalized by
    Grothendieck and later by Fulton and MacPherson.
    \item \textbf{Mirror symmetry:} The HRR theorem underpins
    computations in string theory and the study of Calabi--Yau
    manifolds.
\end{itemize}

\subsection{Topology}

The signature theorem and the theory of multiplicative sequences had a
profound impact on topology:
\begin{itemize}
    \item \textbf{Cobordism theory:} Hirzebruch's genera provided
    invariants that distinguish cobordism classes. Thom's cobordism
    theory combined with Hirzebruch's work to compute the cobordism
    ring.
    \item \textbf{Characteristic classes:} Hirzebruch's systematic
    treatment made characteristic classes a standard tool in topology.
    \item \textbf{Integrality theorems:} The Hirzebruch integrality
    theorems (e.g., the \(\hat{A}\)-genus of a spin manifold is an
    integer) are fundamental in the study of manifolds with special
    holonomy.
    \item \textbf{Index theory:} Hirzebruch's work directly motivated
    the Atiyah--Singer index theorem.
\end{itemize}

\subsection{The Max Planck Institute for Mathematics and the Hirzebruch Prize}

In 1981, Hirzebruch founded the Max Planck Institute for Mathematics
(MPIM) in Bonn. Under his directorship (1981--1994), the institute
became a world leader in mathematics, hosting thousands of visitors and
fostering collaborations across all fields of pure mathematics.

The MPIM established the \textbf{Hirzebruch Prize} in 2018, awarded
biennially to young mathematicians for outstanding contributions to
mathematics. The prize is a fitting tribute to Hirzebruch's vision of
mathematics as a collaborative, international enterprise.

\subsection{Mathematical Physics}

Hirzebruch's work has deep connections to physics:
\begin{itemize}
    \item The HRR theorem is used in string theory to compute the
    number of massless modes in compactifications.
    \item The signature theorem governs anomalies in quantum field
    theory.
    \item Hirzebruch surfaces appear as examples in gauge theory and
    mirror symmetry.
    \item The \(\hat{A}\)-genus appears in the Green--Schwarz anomaly
    cancellation mechanism.
\end{itemize}

\section{Frequently Asked Questions}

\subsection*{Q1: What is the Hirzebruch--Riemann--Roch theorem, and why is it important?}

The Hirzebruch--Riemann--Roch theorem computes the Euler characteristic
\(\chi(X, E) = \sum (-1)^i \dim H^i(X, \mathcal{O}(E))\) of a
holomorphic vector bundle \(E\) on a compact complex manifold \(X\) as
\(\int_X \ch(E) \Td(TX)\). It is important because it reduces a
difficult analytic problem (computing sheaf cohomology dimensions) to a
mechanical topological computation (integrating characteristic
classes). It was the first general Riemann--Roch theorem in higher
dimensions and inspired the Grothendieck--Riemann--Roch theorem and the
Atiyah--Singer index theorem.

\subsection*{Q2: Who was Todd, and why is the Todd class named after him?}

J. A. Todd (1908--1994) was a British algebraic geometer who studied
invariants of algebraic varieties. In the 1930s, he introduced certain
polynomials (now called Todd polynomials) that arise in the study of
Chern classes. Hirzebruch recognized that these polynomials are exactly
the multiplicative sequence associated to \(x/(1-e^{-x})\), and named
the resulting characteristic class the Todd class in Todd's honor.

\subsection*{Q3: How does the signature theorem relate to the index theorem?}

The signature \(\tau(M)\) of a \(4k\)-manifold is the index of a
certain elliptic operator called the signature operator. The Hirzebruch
signature theorem is therefore a special case of the Atiyah--Singer
index theorem. Conversely, Hirzebruch's formulation of the signature as
an integral of the \(L\)-genus provided the model that Atiyah and
Singer generalized to arbitrary elliptic operators.

\subsection*{Q4: What are Hirzebruch surfaces used for?}

Hirzebruch surfaces \(\Sigma_n\) appear in many contexts: they are the
building blocks of rational surfaces in the Enriques--Kodaira
classification; they are the complete smooth toric surfaces; they
provide counterexamples and test cases in the study of algebraic
surfaces; and they appear in symplectic geometry as fundamental examples of
symplectic 4-manifolds with \(b^+ = 1\).

\subsection*{Q5: What is a multiplicative sequence in simple terms?}

A multiplicative sequence is a way to convert a power series
\(Q(x)\) into a sequence of polynomials \(K_n(p_1, \dots, p_n)\) with
the property that if a manifold decomposes as a product, the genus
\(\int_M K(p_1, \dots, p_k)\) factors as the product of the genera of
the factors. The Todd genus and \(L\)-genus are the most famous
examples.

\subsection*{Q6: What open problems are related to Hirzebruch's work?}

Several deep problems remain:
\begin{enumerate}
    \item \textbf{Generalized Riemann--Roch:} The full implications of
    the Grothendieck--Riemann--Roch theorem for motivic cohomology and
    algebraic \(K\)-theory are still being explored.
    \item \textbf{Classification of surfaces:} The complete
    classification of algebraic surfaces of general type involves
    subtle Chern number inequalities that extend Hirzebruch's work.
    \item \textbf{Index theorem for singular spaces:} Generalizing the
    index theorem (and HRR) to singular varieties is an active area.
    \item \textbf{Quantum field theory:} The role of Hirzebruch's genera
    in anomalies, topological field theories, and string theory
    continues to develop.
\end{enumerate}

\section*{Timeline}
\addcontentsline{toc}{section}{Timeline}\[
\begin{array}{l|l}
\textbf{Year} & \textbf{Event} \\ \hline
1927 & \text{Born on October 17 in Breslau, Germany (now Wroc\l aw, Poland)} \\
1945--1947 & \text{Prisoner of war in the Soviet Union} \\
1947--1950 & \text{Studies mathematics at the University of M\"unster} \\
1950 & \text{Ph.D.\ at the University of M\"unster under Heinz Hopf} \\
1952 & \text{Habilitation at the University of M\"unster} \\
1954 & \text{Hirzebruch--Riemann--Roch theorem published} \\
1955 & \text{Visiting the Institute for Advanced Study, Princeton} \\
1956 & \text{Book \textit{Neue topologische Methoden in der
algebraischen Geometrie} appears} \\
1957--1958 & \text{Member of the Institute for Advanced Study, Princeton} \\
1960 & \text{Professor at the University of Bonn} \\
1960s & \text{Hirzebruch surfaces introduced; work on algebraic
surfaces and characteristic classes} \\
1967 & \text{Signature theorem and multiplicative sequences fully developed} \\
1970s & \text{Work on topology of algebraic varieties, the proportionality
principle, and index theory} \\
1981 & \text{Founds the Max Planck Institute for Mathematics in Bonn} \\
1988 & \text{Awarded the Wolf Prize in Mathematics (shared with
Lars H\"ormander)} \\
1990s & \text{Continued work on algebraic geometry and topology} \\
2012 & \text{Passes away on May 27 in Bad Honnef, Germany} \\
2018 & \text{MPIM establishes the Hirzebruch Prize}
\end{array}
\]

\section{Related Chapters}
For further exploration, see Chapter~36 (Bott) on characteristic classes and index theory and Chapter~63 (Donaldson) on 4-manifold topology and gauge theory.

\section*{Exercises for the Reader}
\addcontentsline{toc}{section}{Exercises}

\begin{enumerate}
    \item [B] Let \(C\) be a compact Riemann surface of genus \(g\) and let
    \(L\) be a line bundle of degree \(d\). Use the Hirzebruch--Riemann--
    Roch theorem to compute \(\chi(C, L)\). Show that you recover the
    classical Riemann--Roch theorem by noting that
    \(\chi(C, L) = \ell(L) - \ell(K \otimes L^{-1})\).

    \item [B] Compute the Euler characteristic
    \(\chi(\mathbb{CP}^n, \mathcal{O}(d))\) using HRR and verify the
    formula \(\binom{d+n}{n}\) for \(d \ge 0\). What happens when
    \(d < 0\)?

    \item [B] For the complex projective plane \(\mathbb{CP}^2\), compute
    the Todd class \(\Td(T\mathbb{CP}^2)\) and verify that
    \(\chi(\mathbb{CP}^2, \mathcal{O}) = 1\).

    \item [B] Show that the signature of \(\mathbb{CP}^2\) is \(1\) using
    the Hirzebruch signature theorem. Compute the signature of
    \(\mathbb{CP}^1 \times \mathbb{CP}^1\) using the product formula.

    \item [I] Let \(M\) be a compact oriented \(4\)-manifold with
    \(p_1(M) = 0\). What does the signature theorem say about
    \(\tau(M)\)? Give an example of such a manifold.

    \item [I] Verify that the power series \(Q(x) = x/(1 - e^{-x})\) gives
    the first Todd polynomials:
    \[
    \Td_1 = \frac12 c_1,\quad
    \Td_2 = \frac1{12}(c_1^2 + c_2).
    \]
    (Hint: expand \(\prod_i Q(x_i)\) and express in elementary symmetric
    functions.)

    \item [I] Prove that the Hirzebruch surface \(\Sigma_1\) is isomorphic
    to the blowup of \(\mathbb{P}^2\) at a point. Show that
    \(c_1^2(\Sigma_1) = 8\) and \(c_2(\Sigma_1) = 4\).

    \item [I] For the Hirzebruch surface \(\Sigma_n\), compute the canonical
    divisor \(K_{\Sigma_n}\) and verify the adjunction formula.

    \item [I] Let \(X\) be a compact complex surface. Use HRR to prove the
    Noether formula:
    \[
    \chi(X, \mathcal{O}_X) = \frac1{12}(c_1(X)^2 + c_2(X)).
    \]

    \item [I] Show that the Todd genus of a Calabi--Yau threefold
    (\(c_1 = 0\)) is \(0\). What does HRR give for
    \(\chi(X, \mathcal{O}_X)\)?

    \item [A] Let \(Q(x) = 1 + q_1 x + q_2 x^2 + \cdots\) be a power
    series. Show that the associated multiplicative sequence
    \(\{K_n\}\) satisfies
    \[
    K_n(p_1, \dots, p_n) = \sum_{m_1 + 2m_2 + \cdots + nm_n = n}
    \frac{q_1^{m_1} q_2^{m_2} \cdots q_n^{m_n}}{m_1! m_2! \cdots m_n!}
    p_1^{m_1} p_2^{m_2} \cdots p_n^{m_n}.
    \]

    \item [I] The Hirzebruch--Riemann--Roch theorem can be applied to
    compute the arithmetic genus of a complete intersection. Compute
    \(\chi(X, \mathcal{O}_X)\) for a smooth hypersurface of degree
    \(d\) in \(\mathbb{CP}^n\).

    \item [A] Research the Grothendieck--Riemann--Roch theorem. Explain
    how it generalizes HRR to proper morphisms \(f: X \to Y\) and why
    the term \(\Td(TY)\) appears on the left-hand side.

    \item [B] The Hirzebruch prize was established in 2018. Find the list
    of recipients and describe their research areas. How does the
    prize reflect Hirzebruch's mathematical values?

    \item [A] Let \(M\) be a \(4k\)-dimensional spin manifold. Show that
    the \(\hat{A}\)-genus is an integer. Use the signature theorem and
    the relation between \(L\)-genus and \(\hat{A}\)-genus:
    \[
    L = 2^{2k} \hat{A}(p_1/2, p_2/2, \dots),
    \]
    though the precise relation involves the power series
    \(\sqrt{x}/\sinh\sqrt{x}\).
\end{enumerate}

\section*{Glossary}
\addcontentsline{toc}{section}{Glossary}

\begin{description}
    \item[Arithmetic genus] The Euler characteristic
    \(\chi(X, \mathcal{O}_X) = \sum (-1)^i \dim H^i(X,
    \mathcal{O}_X)\) of the structure sheaf of a complex manifold.
    \item[Bernoulli numbers] Rational numbers \(B_n\) defined by
    \(x/(e^x - 1) = \sum_{n=0}^\infty B_n x^n / n!\).
    \item[Chern character] A ring homomorphism from \(K\)-theory to
    cohomology given by \(\ch(E) = \sum e^{x_i}\) where \(x_i\) are
    Chern roots.
    \item[Chern class] A characteristic class \(c_i(E) \in H^{2i}(X,
    \Z)\) measuring the twisting of a complex vector bundle.
    \item[Chern number] An integer obtained by integrating a product
    of Chern classes over the fundamental class of a manifold.
    \item[Cobordism] Two manifolds are cobordant if their disjoint
    union bounds a third manifold.
    \item[Euler characteristic] The alternating sum of Betti numbers
    \(\chi(M) = \sum (-1)^i \dim H^i(M, \R)\).
    \item[Genus] A multiplicative, additive cobordism invariant of
    manifolds, typically expressed as \(\int_M K(p_1, \dots, p_k)\).
    \item[Grothendieck--Riemann--Roch] A generalization of HRR to
    proper morphisms, relating the pushforward of the Chern character
    to the Chern character of the direct image.
    \item[Hirzebruch surface] The projective bundle
    \(\Sigma_n = \mathbb{P}(\mathcal{O} \oplus \mathcal{O}(-n))\) over
    \(\mathbb{P}^1\).
    \item[Hirzebruch--Riemann--Roch theorem] The formula
    \(\chi(X, E) = \int_X \ch(E) \Td(TX)\) for a holomorphic vector
    bundle on a compact complex manifold.
    \item[Hirzebruch signature theorem] The formula
    \(\tau(M) = \int_M L(p_1, \dots, p_k)\) for a \(4k\)-manifold.
    \item[Index theorem] The Atiyah--Singer theorem: the index of an
    elliptic operator equals a topological expression in
    characteristic classes.
    \item[\(K\)-theory] The Grothendieck group of vector bundles over
    a space, with addition given by direct sum.
    \item[\(L\)-genus] The multiplicative sequence associated to
    \(\sqrt{x}/\tanh\sqrt{x}\), giving the signature of a manifold.
    \item[Multiplicative sequence] A sequence of polynomials
    \(K_n(p_1, \dots, p_n)\) satisfying \(K(p+q) = K(p) K(q)\) under
    the natural product of characteristic classes.
    \item[Noether formula] The relation
    \(\chi(S, \mathcal{O}_S) = (c_1^2 + c_2)/12\) for a complex
    surface.
    \item[Pontryagin class] A characteristic class
    \(p_i(M) \in H^{4i}(M, \Z)\) of an oriented real vector bundle,
    obtained from Chern classes of the complexification.
    \item[Pontryagin number] An integer obtained by integrating a
    product of Pontryagin classes over the fundamental class of an
    oriented manifold.
    \item[Proportionality principle] The principle that Chern numbers
    of compact locally symmetric spaces are proportional to those of
    the compact dual.
    \item[Signature] The number of positive minus negative eigenvalues
    of the intersection form on the middle cohomology of a
    \(4k\)-manifold.
    \item[Splitting principle] The principle that for computations
    with characteristic classes, one may assume a vector bundle
    splits as a sum of line bundles.
    \item[Todd class] The multiplicative sequence associated to
    \(x/(1 - e^{-x})\), used in the HRR theorem.
    \item[Todd genus] The integer \(\int_X \Td(TX)\) for a complex
    manifold \(X\), equal to the arithmetic genus.
    \item[\(\hat{A}\)-genus] The multiplicative sequence associated to
    \(\sqrt{x}/\sinh\sqrt{x}\), giving the index of the Dirac operator
    on a spin manifold.
\end{description}

\section{Citations}

\subsection*{Primary Sources}
\begin{enumerate}
    \item Hirzebruch, F. (1956). \textit{Neue topologische Methoden in der algebraischen Geometrie}. Ergebnisse der Mathematik und ihrer Grenzgebiete, N.F., Heft 9. Springer-Verlag.
    \item Atiyah, M. F., \& Hirzebruch, F. (1959). Riemann-Roch theorems for differentiable manifolds. \textit{Bulletin of the American Mathematical Society}, 65(4), 276--281.
    \item Atiyah, M. F., \& Hirzebruch, F. (1961). Vector bundles and homogeneous spaces. In \textit{Proceedings of Symposia in Pure Mathematics}, Vol. 3, 7--38. American Mathematical Society.
    \item Hirzebruch, F. (1966). \textit{Topological Methods in Algebraic Geometry}. 3rd ed. Grundlehren der mathematischen Wissenschaften, Vol. 131. Springer-Verlag.
\end{enumerate}

\subsection*{Biographies and Overviews}
\begin{enumerate}
    \item Atiyah, M. F. (2012). Friedrich Hirzebruch (1927--2012). \textit{Notices of the American Mathematical Society}, 59(10), 1374--1377.
    \item Zagier, D. (2012). Friedrich Hirzebruch (1927--2012). \textit{Jahresbericht der Deutschen Mathematiker-Vereinigung}, 114(3), 131--140.
    \item Eichler, M., \& Zagier, D. (Eds.). (1987). \textit{Arbeitstagung Bonn 1987: Proceedings of the Arbeitstagung held in Bonn, June 15-21, 1987}. Progress in Mathematics, Vol. 75. Birkhäuser.
\end{enumerate}

\subsection*{Textbooks and Expository Works}
\begin{enumerate}
    \item Milnor, J. W., \& Stasheff, J. D. (1974). \textit{Characteristic Classes}. Annals of Mathematics Studies, Vol. 76. Princeton University Press.
    \item Bott, R., \& Tu, L. W. (1982). \textit{Differential Forms in Algebraic Topology}. Graduate Texts in Mathematics, Vol. 82. Springer-Verlag.
    \item Palais, R. S. (1965). \textit{Seminar on the Atiyah-Singer Index Theorem}. Annals of Mathematics Studies, Vol. 57. Princeton University Press.
    \item Hatcher, A. (2009). \textit{Vector Bundles and K-Theory}. Available online at \url{https://pi.math.cornell.edu/~hatcher/VBKT/VBpage.html}.
\end{enumerate}
